\section{Introduction}
\label{sec:intro}
Monitoring a newly introduced target presents an asymmetric learning problem. The target can supply examples of legitimate operation, while other clients may already have labeled malicious traffic. A supervised detector trained only for the target is unavailable, but discarding peer attack information would leave the monitor dependent on deviations from its own benign sample. The practical question is how to preserve useful peer evidence while retaining a local route for attacks that those peers miss.

A common fusion strategy takes the largest peer score and calibrates one threshold on target benign observations. This strategy can lose discrimination when the benign tail of one peer dominates the fused score: the common threshold must accommodate that tail, and another peer's informative but numerically smaller attack evidence may no longer trigger. We call this task-specific failure mode peer drowning. Score normalization and domination by a single extreme detector in maximum fusion are established concerns in outlier ensembles~\cite{zimek2014}. Target-relative score transformations address differences in score scale, but they do not make the tails of heterogeneous experts interchangeable or ensure that a single threshold on their maximum preserves every useful detector.

PACE, Peer-wise Calibrated Experts, retains the supervised peers as separate experts. Each margin is mapped relative to scores on a target benign reference role. Separate empirical thresholds are calibrated on a later benign role, and a shared multiplier adjusts the peer union according to its observed blockwise alert rate. A protected local path combines a six-detector mean anchor with feature-rarity and peer-response experts. Its purpose is to detect attacks outside the peer union, rather than merely improving the detection rate of a local model on attacks already covered by peers.

The components are deliberately conventional. PACE does not introduce a new isolation estimator, density model, neural architecture, or multiple-testing theorem. Moreover, calibrated fusion of supervised and benign-deviation evidence has close prior art, including MCDE~\cite{yue2026mcde}. The question investigated here is narrower: whether retaining foreign peers separately and replacing a mean-ensemble local branch with an anchored heterogeneous bank improves this particular target-label-free benchmark setting. The distinction is architectural and empirical; a claim of being the first general calibrated hybrid intrusion detector would not be justified.

Three aspects of the evaluation constrain that question. First, the fitting budget counts rows used by local detectors, while additional target benign data contribute to preprocessing, reference construction, and calibration. Second, the nominal operating parameter does not provide a test-time false-alert guarantee. The implementation's resolution rule can even allocate more than that parameter across the three decision paths. Third, the three reported seeds measure learning variability within the same benchmark cohort. Conclusions are therefore specific to this evaluation setting.

The article makes three concrete contributions. It specifies the implemented peer-wise and local-tier decision rule, including its finite-sample quantile conventions and allocation accounting. It evaluates the local-branch replacement on 58 targets from six corpora and seven study strata, using paired primary comparisons, false-alert distributions, ablations, and within-cohort budget changes. Finally, it identifies the boundaries of the observed benefit: concentration in TON\_IoT, limited gains in several other strata, a CTU-13 target regression, and weaker results at larger local budgets.

At nominal $\alpha=0.01$, PACE improves target-macro detection over the mean-ensemble local comparator by 0.026086 with 100 local fitting rows and by 0.028306 with 400 rows. These gains coexist with mean realized false-alert rates of 0.033992 and 0.032369 and worst-target rates close to 0.67. The appropriate conclusion is that complementary local experts can improve detection in this offline setting, while calibration and deployment validity remain material limitations. The study does not establish an operational one-percent false-alert detector.

